\chapter{Discussione e conclusioni}
\label{cap:conclusioni}

% ---------------------------------------------------------------------
% Capitolo aggiunto: nella struttura precedente le conclusioni erano il
% Cap. 6, ora occupato dalla validazione sul campo.
%
% Scritto il 2026-08-17. I contributi sono presentati come METODOLOGICI:
% il valore del lavoro non e' un F2 piu' alto di altri, ed e' importante
% che il capitolo non provi a farlo sembrare tale.
%
% NON usare il "4,32x sull'escursione dell'angolo del tronco": ISS-033 ha
% accertato che non si riproduce (rieseguendo su 25 sequenze esce 1,99x).
% Il numero da citare e' il 2,9x sull'accuratezza, oppure il divario
% 0,438 in F2 fra le due camere, che sono misure indipendenti su 5 fold.
% ---------------------------------------------------------------------

\section{Contributi}
\label{sec:contributi}

Il lavoro non rivendica un record di accuratezza, e sarebbe poco onesto farlo: i numeri pubblicati su UP-Fall sono più alti di quelli riportati qui. Rivendica invece che quei numeri e questi rispondano a domande diverse, e che la differenza sia misurabile.

\paragraph{1. Un protocollo di valutazione per dataset con pochi soggetti, e la misura di cosa succede senza.} Con dieci soggetti, una divisione casuale delle finestre sovrapposte mette tutte e dieci le persone sia in addestramento sia in verifica. Applicando insieme le tre correzioni --- divisione per soggetto, normalizzazione delle scale, soglia scelta su un terzo insieme --- alla stessa configurazione e sugli stessi dati, l'F1 passa da $0{,}0000$ a $0{,}234$. Il primo valore non è un modello mediocre: è un modello che non ha mai predetto la classe positiva, e che senza il protocollo corretto sarebbe stato indistinguibile da uno funzionante.

Il contesto rende il contributo leggibile. I lavori pubblicati su UP-Fall riportano accuratezze fra il 98\% e il 99,58\%, e almeno uno dichiara esplicitamente di usare partizioni casuali 70:30 dell'intero dataset --- cioè esattamente il difetto qui documentato. Un riscontro indipendente quantifica lo scarto: lo stesso classificatore che dà 89\% in \emph{k-fold} casuale scende al 76\% con divisione per soggetto.

\paragraph{2. La quantificazione dell'effetto del punto di vista.} \emph{Misura condotta sull'impianto a cinque soggetti, non ripetuta su dieci; l'ampiezza dell'effetto è tale che il riallineamento non può ribaltarla, i valori assoluti vanno letti come indicativi.} Il passaggio dalla telecamera laterale a quella frontale rialzata, a parità di soggetto, di caduta e di istante, costa $0{,}438$ di F2 --- da $0{,}872$ a $0{,}434$. La ragione è geometrica e non algoritmica: una telecamera che guarda dall'alto comprime la rotazione del tronco lungo il proprio asse ottico, e una persona eretta e una distesa si proiettano in modo simile. L'informazione che distingue una caduta viene cancellata dall'ottica prima che un modello possa elaborarla. Addestrando su entrambi i punti di vista se ne recupera il 56\%, con i falsi positivi in calo del 59\%: migliorano insieme, il che distingue un apprendimento più generale da uno spostamento di soglia.

\paragraph{3. Un confronto fra classificatori a parità di protocollo, con esito contrario all'impostazione di partenza.} Cinque architetture --- due ricorrenti, una convoluzionale, un encoder ad attenzione e un classificatore non sequenziale --- cadono entro $0{,}066$ di $F_2$, contro una dispersione fra soggetti compresa fra $0{,}08$ e $0{,}12$. Su dieci soggetti \textbf{nessuna coppia è distinguibile}, e il numero di parametri non ordina il risultato.

Il contributo non è però la constatazione, che da sola sarebbe un risultato negativo: è la sua quantificazione. Il confronto appaiato permette di calcolare quante persone servirebbero per rendere significativa ciascuna differenza, e la risposta separa le domande in due classi. Distinguere una GRU da una LSTM richiederebbe \textbf{centotredici} soggetti; distinguere due definizioni del compito ne richiede \textbf{due}. Un fattore cinquantasei fra le due domande, e nessun archivio pubblico di cadute si avvicina alla prima cifra.

\paragraph{4. La caratterizzazione end-to-end su hardware edge reale.} L'intera catena gira su un dispositivo da circa cento euro a 26,5 fotogrammi al secondo, senza che alcun dato lasci l'abitazione. La ripartizione del costo per fotogramma è il dato meno atteso: la percezione è l'\textbf{88,4\%} del tempo, il classificatore il \textbf{5,8\%}. Tutta la letteratura sul confronto fra architetture, e tutto il lavoro del punto 3, agiscono su un ventesimo del problema.

\paragraph{5. La quantificazione dello scarto fra precision di laboratorio e falsi allarmi in esercizio.} Una precision di $0{,}819$ misurata su un insieme quasi bilanciato corrisponde, proiettata sul funzionamento continuo, a 14--125 falsi allarmi all'ora. Misurati: \textbf{2.735}. La previsione, costruita per essere scomoda, era ottimistica di oltre un ordine di grandezza.

\paragraph{6. Le grandezze integrali al posto delle derivate istantanee.} Il classificatore non usava la dinamica --- correlazione fra punteggio e movimento $+0{,}139$ --- non per un difetto di architettura ma perché ciò che riceveva come dinamica era rumore: le derivate hanno autocorrelazione a un fotogramma fra $0{,}11$ e $0{,}41$, contro $0{,}84$--$0{,}93$ delle grandezze statiche. Sostituendole con grandezze accumulate su una finestra mobile, l'autocorrelazione sale a $0{,}91$--$0{,}94$ e la separazione fra le classi raggiunge fattori di 21--32 nelle mediane. È l'unico intervento del lavoro che migliora simultaneamente in laboratorio e sul campo.

\paragraph{7. La misura di ciò che il dataset di riferimento non riesce a distinguere.} Fra la configurazione di riferimento e quella adottata corrono $0{,}055$ di $F_2$ in laboratorio --- dentro la dispersione fra soggetti, quindi una differenza che non si può dichiarare --- e un fattore \textbf{cinquantadue} sui falsi allarmi in una casa. È il contributo che si ritiene principale, ed è discusso nella sezione seguente.

\section{Il risultato che non era atteso}
\label{sec:risultato-inatteso}

Il lavoro era impostato su una domanda di ingegneria: \emph{quale architettura, e su quale hardware}. La risposta è che la domanda era mal posta, e per due ragioni indipendenti.

La prima è che \textbf{l'architettura del classificatore non è la variabile che decide}. L'intero intervallo entro cui cadono cinque architetture vale $0{,}066$; il divario fra due posizioni della telecamera vale $0{,}438$, cioè \textbf{quasi sette volte tanto}, e quello fra due definizioni del compito $0{,}127$ con una differenza statisticamente netta, mentre nessuna coppia di architetture lo è. Ha una conseguenza pratica immediata per chi installa questi sistemi, ed è economicamente rilevante: spostare una telecamera costa meno che addestrare un modello più grande, e rende di più.

La seconda, e più importante, riguarda \textbf{cosa un insieme di verifica può
misurare}. Tre interventi progettati per la robustezza domestica sono stati
giudicati da un archivio di laboratorio e da una casa reale, con gli stessi
modelli e la stessa procedura. Fra la configurazione di riferimento e quella
adottata corrono $0{,}055$ di $F_2$ in laboratorio --- \emph{dentro} la
dispersione fra soggetti, quindi una differenza che il protocollo non
autorizza a dichiarare --- e un fattore \textbf{cinquantadue} sulle finestre in
allarme misurate in una casa.

Lo stesso fenomeno si ripresenta su un singolo iperparametro: portando la
finestra di accumulazione da 31 a 45 fotogrammi l'$F_2$ di laboratorio
continua a salire mentre sul campo si perde una caduta.

\paragraph{Una precisazione che il lavoro deve a sé stesso.} Una versione
precedente di questa tesi sosteneva qualcosa di più forte: che l'archivio di
laboratorio \emph{ordinasse in modo errato} le alternative, collocando ultima
la configurazione che in casa risultava migliore. Quella affermazione era
fondata su cinque soggetti, e non è sopravvissuta al raddoppio: su dieci, le
due graduatorie concordano.

L'enunciato che regge è quindi più debole in apparenza e più solido in
sostanza:

\begin{quote}
Un insieme di verifica non può misurare un miglioramento su una condizione che
non rappresenta. Non si limita a sovrastimare le prestazioni: \textbf{rende
indistinguibili alternative che in esercizio differiscono di ordini di
grandezza}, e lo fa con la stessa apparente sicurezza con cui distinguerebbe
alternative davvero equivalenti.
\end{quote}

L'effetto pratico coincide con quello di un ordinamento errato. Chi selezioni
la configurazione sul solo $F_2$ di laboratorio non ha alcuna ragione
statistica per preferire la configurazione adottata a quella di riferimento:
sono indistinguibili, e la seconda è più semplice. Sceglierebbe la più
semplice, e installerebbe un sistema che produce cinquantadue volte i falsi
allarmi --- avendo seguito la procedura corretta fino in fondo.

La causa non è misteriosa ed è stata misurata. Nell'archivio di laboratorio le
persone orizzontali che non sono cadute stanno \emph{a terra}, perché
sdraiarsi sul pavimento è una delle attività registrate. In una casa stanno
\emph{sul divano}. Un intervento che insegni a distinguere il divano dal
pavimento non può essere premiato da un insieme che non contiene divani ---
e vi consuma anzi capacità del modello su una distinzione che lì non serve.

La differenza fra le due formulazioni è operativa. Che un dataset sovrastimi è noto, ed è misurato su scala ben maggiore di questo lavoro: il benchmark OmniFall, unificando otto dataset pubblici e confrontandoli con cadute reali, riporta per UP-Fall un crollo da 72,9\% a 16,9\% di F1. Che un dataset \emph{ordini male} è una proprietà diversa, e più dannosa: chi sovrastima produce un sistema peggiore di quanto crede; chi non distingue produce \emph{il sistema sbagliato}, avendo scartato quello giusto con piena fiducia nella misura e senza che nulla, nella procedura, potesse segnalarglielo.

\section{Limiti}
\label{sec:limiti}

Si dichiarano senza attenuazioni, e alcuni sono sufficienti a delimitare quanto sopra si può portare.

\begin{itemize}
  \item \textbf{Dieci soggetti} nel dataset di addestramento, tutti giovani adulti, contro una popolazione bersaglio anziana. La dispersione fra fold, $\pm 0{,}05$--$0{,}14$, che il raddoppio dei soggetti non ha ridotto in modo apprezzabile, è dell'ordine di molte differenze confrontate nel Capitolo~\ref{cap:risultati}, che vanno quindi dichiarate non distinguibili.

  \item \textbf{Le cadute sono recitate}, qui come in UP-Fall, in UR-Fall e nella quasi totalità della letteratura. Nessuno studio di pochi giorni può misurare cadute accidentali: una persona anziana cade meno di una volta all'anno. È il limite che questo lavoro condivide con il settore, e va detto invece che lasciato intendere.

  \item \textbf{Il recall sul campo poggia su cinque cadute}, in una sessione con l'inquadratura difettosa. Non è una stima, è un indizio. E non è mai misurato sulla popolazione anziana, che per ragioni di sicurezza non ha eseguito cadute.

  \item \textbf{Un solo archivio esterno} per la validazione fra dataset: non si distingue ciò che è specifico di UR-Fall da ciò che varrebbe per qualunque ambiente nuovo.

  \item \textbf{Due sole abitazioni} per i controesempi domestici. Parte del guadagno misurato potrebbe essere adattamento a quelle due stanze e non robustezza generale, e servirebbe una terza casa per separare le due cose.

  \item \textbf{Nessuna ricerca sistematica di iperparametri}, in nessun esperimento. Differenze inferiori a $\sim 0{,}02$ fra modelli neurali non sono interpretabili.

  \item \textbf{Il tracciamento multi-persona non è validato}, e non lo è per assenza di dati: nessun archivio pubblico contiene più persone in scena.

  \item \textbf{L'effetto della quantizzazione INT8 non è stato isolato} dagli altri fattori nel confronto fra stimatori di posa.

  \item \textbf{Le occlusioni da mobilio non sono state misurate.} Il risultato sull'auto-occlusione riguarda YOLOv8, e rifatto sullo stimatore effettivamente installato non si riproduce: su MediaPipe la confidenza del volto passa da $0{,}997$ a $0{,}989$ fra la posizione eretta e quella a terra, cioè non degrada. Degradano gli arti. Il risultato originale va attribuito allo stimatore e non al problema.

  \item \textbf{La prova di durata da dodici ore non è stata eseguita.} La sospensione automatica del dispositivo, che aveva reso invalida senza segnalarlo la prima sessione lunga, è ora rilevabile, ma la prova resta da fare.
\end{itemize}

\section{Sviluppi}
\label{sec:sviluppi}

\paragraph{Taratura della soglia senza etichette, sul tasso di base stimato.} Il crollo cross-dataset si scompone in due parti: due terzi sono taratura della soglia --- il modello continua a ordinare correttamente le finestre --- e un terzo è perdita reale. La parte di taratura si correggerebbe con cadute etichettate dell'ambiente di destinazione, che chi installa non possiede. Fissare la soglia a un percentile della distribuzione dei punteggi osservati recupera il 70\% di quel divario su UR-Fall, portando l'F2 da $0{,}073$ a $0{,}449$; ma richiede di assumere il tasso di base, e su una casa quell'assunzione produrrebbe fra 84 e 1.690 falsi allarmi all'ora, perché fissa il tasso di allarme invece di misurarlo. La direzione promettente non è assumere il tasso ma \emph{stimarlo}, per esempio da un periodo iniziale in cui si assuma l'assenza di cadute --- ipotesi realistica in una casa e falsa in un archivio.

\paragraph{Mappa del contesto tramite rilevamento degli oggetti.} La distinzione fra divano e pavimento è oggi affidata a una stima automatica del piano del pavimento. Un rilevamento dei mobili eseguito \emph{all'avvio} e non per fotogramma --- i mobili non si muovono --- avrebbe costo nullo in esercizio e fornirebbe il contesto in modo diretto.

\paragraph{Sessioni più lunghe e in più abitazioni}, con partecipanti anziani per il solo blocco di vita normale, e la prova di durata da dodici ore.

\paragraph{Sintesi di punti di vista per riproiezione dello scheletro}, per attaccare il divario più grande misurato in questo lavoro. Va però dichiarato un limite noto del metodo: ruotare uno scheletro produce le coordinate di una vista nuova ma non la degradazione di confidenza che quella vista causerebbe realmente, e con auto-occlusioni in direzioni diverse lo scheletro trasformato non imita fedelmente il nuovo punto di vista. Il guadagno riportato in letteratura va quindi preso come limite superiore.

\paragraph{Ensemble di architetture.} I benchmark cross-dataset più recenti indicano negli ensemble il guadagno di generalizzazione più affidabile fra le strategie confrontate. Cinque modelli addestrati a parità di protocollo sono già disponibili, quindi il costo dell'esperimento è quasi nullo.

\paragraph{Ricerca di iperparametri}, per rendere interpretabili le differenze fra modelli neurali che questo lavoro dichiara non distinguibili.

\section{Conclusione}

L'ipotesi economica di partenza è verificata: un dispositivo da circa cento euro esegue in tempo reale l'intera catena di rilevamento, senza che alcuna immagine lasci l'abitazione, e con il classificatore che occupa il 5,8\% del budget computazionale.

L'ipotesi tecnica di partenza --- che la scelta dell'architettura ricorrente fosse la variabile determinante --- non è verificata, e la sua confutazione è il risultato più utile del lavoro. La variabile che decide è dove si mette la telecamera, e il fattore che più di ogni altro distorce le decisioni progettuali è la distanza fra l'insieme su cui si misura e l'ambiente in cui si installa.

Il sistema, allo stato attuale, non è pronto per l'installazione: fra i suoi 14 falsi allarmi all'ora e il criterio clinico di uno o due alla settimana restano due o tre ordini di grandezza. Ma la distanza è ora misurata invece che ignorata, il suo meccanismo è spiegato, ed è stato mostrato che l'insieme di verifica su cui il settore misura non solo non la rivela, ma indica la direzione opposta a quella che la riduce.
